\documentclass[11pt]{article}
% Shared typesetting for ECO 202 papers; contains formatting, not answers.
\usepackage[T1]{fontenc}
\usepackage{lmodern}
\usepackage[margin=0.85in]{geometry}
\usepackage{amsmath,amssymb,booktabs,array,enumitem,fancyhdr,lastpage,needspace}
\usepackage[hidelinks]{hyperref}
\setlength{\parindent}{0pt}
\setlength{\parskip}{5pt}
\setlength{\headheight}{14pt}
\setlength{\emergencystretch}{2em}
\pagestyle{fancy}
\fancyhf{}
\lhead{ECO 202 \textbar\ Fall 2026}
\rhead{\paperlabel}
\cfoot{\thepage\ of \pageref{LastPage}}
\newcommand{\paperlabel}{Statistics and Data Analysis for Economics}
% Arguments: complete paper title, date/term, covered class numbers.
\newcommand{\examheader}[3]{%
  \renewcommand{\paperlabel}{#1}%
  \begin{center}
  {\Large\bfseries ECO 202 --- #1}\\[4pt]
  Statistics and Data Analysis for Economics \textbar\ Princeton University\\[3pt]
  #2 \textbar\ Classes #3 \textbar\ 100 points
  \end{center}
  Name: \rule{2.6in}{0.4pt}\hfill Student ID: \rule{1.3in}{0.4pt}\par
  \begin{small}
  \textbf{Timing.} Target workload: 70 minutes (10 minutes multiple choice; 60 minutes written work). Follow announced start/end times within the 80-minute meeting. For practice, set a 70-minute timer.

  \textbf{Conditions.} No books, notes, calculators, AI, or other electronics; University-authorized accommodations apply. Fractions, radicals, and equivalent exact expressions are acceptable. Show reasoning in written answers. Data are teaching examples unless identified otherwise.

  \textbf{Points.} Questions 1--4: one answer A--D, 5 points each, no negative marking. Questions 5--8: 20 points each; subpart points are shown. Label any continuation clearly.
  \end{small}
  \section*{Multiple choice \normalfont\small(20 points; about 10 minutes)}
}
\newcounter{mcquestion}
\newcommand{\mcq}[5]{%
  \par\addvspace{6pt}\begin{minipage}{\linewidth}
  \refstepcounter{mcquestion}\textbf{\themcquestion.} #1
  \begin{enumerate}[label=\textbf{\Alph*.},leftmargin=1.9em,itemsep=1pt,topsep=3pt]
  \item #2
  \item #3
  \item #4
  \item #5
  \end{enumerate}
  \end{minipage}\par
}
\newcommand{\written}[2]{\clearpage\section*{#1. #2 \normalfont\large(20 points; about 15 minutes)}}
\newlist{parts}{enumerate}{1}
\setlist[parts]{label=\textbf{(\alph*)},leftmargin=2em,itemsep=7pt,topsep=5pt}
\newcommand{\pts}[1]{\textbf{(#1 points)}\ }
\newcommand{\workspace}[1]{\par\vspace{#1}}
% Arguments: complete paper title, covered class numbers.
\newcommand{\solutionheader}[2]{%
  \renewcommand{\paperlabel}{#1 solutions}%
  \begin{center}
  {\Large\bfseries ECO 202 --- #1}\\[4pt]
  {\large Worked solutions}\\[3pt]
  Fall 2026 \textbar\ Classes #2
  \end{center}
  Equivalent exact expressions and valid alternative reasoning are acceptable. For practice study, attempt the paper first, then compare reasoning, close these solutions, and reconstruct any missed method independently.
}
\newcommand{\solution}[2]{\Needspace{5\baselineskip}\section*{#1. #2}}

\begin{document}
\examheader{Practice Exam 1 --- Version B}{Fall 2026}{1--5}
\mcq{A dataset records the three-digit telephone area code of each customer. Which description fits area code?}
{Quantitative, because it is written using digits.}
{A measure of the amount of telephone service.}
{A continuous outcome.}
{Categorical, because the numbers label groups.}
\mcq{Under a continuous density model, the fraction of values between $a$ and $b$ is represented by:}
{the sum of the density heights at $a$ and $b$.}
{the area under the density between $a$ and $b$.}
{the density height at the midpoint.}
{the interval length, regardless of the density.}
\mcq{A regression predicts weekly sales in thousands of dollars from store area in hundreds of square feet. Which describes the slope's units?}
{Thousands of dollars per hundred square feet.}
{Hundreds of square feet per thousand dollars.}
{No units; every regression coefficient is standardized.}
{Thousands of dollars squared.}
\mcq{For a treated worker with an observed outcome $Y_i(1)$, the missing counterfactual is:}
{the same worker's outcome under no treatment, $Y_i(0)$.}
{the observed mean of the treated workers.}
{the outcome of any untreated worker.}
{the worker's outcome measured after treatment.}

\written{5}{A boxplot and a change of scale}
Six values are $1,3,3,5,7,8$. Use the median-of-halves quartile convention. For variance you may use $\sum x_i^2=157$ and $s^2=[\sum x_i^2-n\bar x^2]/(n-1)$.
\begin{parts}
\item \pts{6} Calculate the median, quartiles, and IQR. Sketch a labeled boxplot with whiskers extending to the extreme observations within the $1.5\mathrm{IQR}$ fences.
\workspace{1.1in}
\item \pts{6} Calculate the mean and sample variance. State how the units of variance differ from the units of the observations.
\workspace{0.85in}
\item \pts{4} Transform each value by $w=20-2x$. Find the transformed median and IQR without listing every new value. Explain how the negative multiplier changes the ordering and why the IQR remains nonnegative.
\workspace{0.8in}
\item \pts{4} Another dataset has $Q_1=3$, $Q_3=7$, and an observation of 20. Apply the outlier rule and explain what you would need to check before deleting that observation.
\workspace{0.8in}
\end{parts}

\written{6}{Transformations and model-based comparisons}
\begin{parts}
\item \pts{6} A density has constant height $1/4$ on $[-2,2]$ and is zero elsewhere. Sketch it and find its median. What fraction of its values satisfies $(X+2)/4<3/4$? Show the transformed cutoff and area.
\workspace{1.1in}
\item \pts{5} In a separate teaching model, delivery time $T$ follows $\mathsf N(30,5^2)$ in minutes. Find the fraction below 20 minutes using $\Phi(-2)=0.0228$. Explain why a model standard deviation of 25 would be incorrect here.
\workspace{0.8in}
\item \pts{5} Company A's delivery times are Normal with mean 30 and standard deviation 5; Company B's are Normal with mean 40 and standard deviation 10. Compare a 35-minute delivery at A with a 45-minute delivery at B, first in minutes and then in percentile position. Use $\Phi(1)=0.8413$ and $\Phi(0.5)=0.6915$.
\workspace{0.9in}
\item \pts{4} An analyst claims that 95\% of every dataset must lie within roughly two standard deviations of its mean. Explain the problem and identify one graphical feature that would lead you to question a Normal model.
\workspace{0.85in}
\end{parts}

\written{7}{A small scatterplot and its residuals}
Three observed pairs are $(x,y)=(0,1),(1,3),(2,2)$, with $x$ measured in hours and $y$ in units of output. You may use
\[S_{xx}=\sum(x_i-\bar x)^2,\quad S_{yy}=\sum(y_i-\bar y)^2,\quad S_{xy}=\sum(x_i-\bar x)(y_i-\bar y),\]
\[b_1=S_{xy}/S_{xx},\qquad b_0=\bar y-b_1\bar x,\qquad r=S_{xy}/\sqrt{S_{xx}S_{yy}}.\]
\begin{parts}
\item \pts{6} Find the two means, the three centered sums, and the least-squares line with an intercept.
\workspace{1.0in}
\item \pts{4} Sketch the three observations and fitted line on labeled axes. Mark the vertical residual for the point $(1,3)$ and calculate it.
\workspace{1.15in}
\item \pts{6} Calculate $r$, $r^2$, and all three residuals. Check that residuals sum to zero. Explain what $r^2$ measures.
\workspace{0.85in}
\item \pts{4} A fourth observation with $x=50$ is added. Explain why it has high leverage and how you would assess whether it is influential. Can its effect on the fitted slope be determined from its $x$ value alone?
\workspace{0.8in}
\end{parts}

\written{8}{Group composition and a causal claim}
The following fictional observational records describe training participants and nonparticipants. Prior skill was measured before training, and all outcomes are subsequent earnings in thousands of dollars.
\begin{center}\begin{tabular}{lrrrr}\toprule
&\multicolumn{2}{c}{Training}&\multicolumn{2}{c}{No training}\\
Prior skill & Count & Mean earnings & Count & Mean earnings\\\midrule
High &1&12&3&10\\Low&3&6&1&4\\\bottomrule\end{tabular}\end{center}
\begin{parts}
\item \pts{6} Calculate mean earnings among all participants and among all nonparticipants, and their difference. Show how the counts enter the calculation.
\workspace{0.95in}
\item \pts{4} Calculate the participant-minus-nonparticipant mean difference separately within each prior-skill group. Compare the signs with part (a).
\workspace{0.7in}
\item \pts{6} Explain how the group composition produces the contrast between (a) and (b). Does either comparison by itself establish a causal training effect? Explain what remains uncertain even after separating workers by prior skill.
\workspace{1.1in}
\item \pts{4} For a worker who participated, define the missing counterfactual relevant to this question. Explain why another worker in the same prior-skill group need not supply that counterfactual exactly.
\workspace{0.85in}
\end{parts}
\end{document}
